\chapter{Conclusiones y trabajos futuros}
\label{chap:conclusiones}

Este capítulo cierra la memoria valorando con cuidado en qué medida se han alcanzado los objetivos planteados, qué conocimientos del grado se han puesto en práctica, qué se ha aprendido en el proceso y qué líneas quedan abiertas. La idea que ha vertebrado todo el trabajo vuelve aquí como criterio de evaluación: no bastaba con construir un RAG jurídico que «funcione», sino con \emph{medir con evidencia} qué decisiones lo hacen fiable y reconocer, sin adornos, dónde la evidencia todavía no alcanza.

\section{Consecución de objetivos}
\label{sec:consecucion-objetivos}

El objetivo general —diseñar, implementar y evaluar de forma rigurosa un sistema RAG sobre legislación consolidada del BOE, poniendo el énfasis en medir qué decisiones producen respuestas más fiables— se considera alcanzado en su parte central y honestamente delimitado en lo que quedó fuera. El repaso a los objetivos específicos lo concreta.

\begin{description}
  \item[OE-01 y OE-02 (corpus y representación).] \textbf{Conseguidos.} Se construyó un corpus reproducible de 92 normas consolidadas, con ingesta cruda inmutable verificada por sumas de comprobación, y una representación \emph{parent-child} que preserva la estructura jurídica y resuelve la vigencia temporal con una regla estricta y una cuarentena prudente para lo irresoluble.

  \item[OE-03 (estrategias de recuperación).] \textbf{Conseguido.} Se implementaron y evaluaron las recuperaciones léxica (BM25 con analizador de español), semántica (índice denso exacto) e híbrida (fusión por rangos recíprocos y combinación convexa). La \emph{reordenación} mediante \emph{cross-encoder} queda como \textbf{trabajo futuro}, con una base empírica que ayuda a acotar su prioridad: la ablación de fusión mostró que añadir la señal léxica no mejora al denso, de modo que un reordenador léxico-informado es poco prometedor; y la descomposición por \emph{baselines} situó el margen de la recuperación no en el orden de los artículos —donde el denso está casi al techo— sino en el ensamblado del contexto, que un \emph{reranker} de artículos no aborda. Un \emph{cross-encoder}, que aprende interacciones consulta--documento distintas de la mera fusión de señales, podría no obstante mejorar el orden; es una línea abierta que este resultado no cierra.

  \item[OE-04 (comparación de recuperadores).] \textbf{Conseguido; es el resultado central.} Sobre un banco anotado mediante \emph{pooling} diversificado para reducir el sesgo de selección, con protocolo estadístico pareado y corrección por comparaciones múltiples, se mostró que \textbf{ninguna variante híbrida supera al recuperador denso, sin que se detecte tampoco una diferencia concluyente frente al mejor híbrido}, mientras que el léxico queda muy por detrás, con especial claridad en las consultas ciudadanas. La conclusión defendible es que \emph{la fusión no aporta} sobre este corpus; el resultado refuta además un hallazgo preliminar, obtenido sobre un corpus reducido, en el que la fusión sí parecía ayudar.

  \item[OE-05 (generación fundamentada y prudente).] \textbf{Conseguido.} El generador redacta a partir de las evidencias, cita solo identificadores entregados y está diseñado para abstenerse cuando no detecta evidencia suficiente, aunque no lo consigue en todos los casos: no produjo respuestas indebidas en las 30 preguntas de dominio lejano, pero sí respondió indebidamente en una de las 10 preguntas \emph{near-miss}. Se analiza además su calibración conservadora (una sobre-abstención apreciable en preguntas respondibles).

  \item[OE-06 (marco de evaluación y validación del juez).] \textbf{Conseguido, con un resultado negativo.} Se diseñó un marco modular de seis capas y se validó el modelo juez frente a anotación humana. La validación —que era justamente lo que exigía el objetivo— reveló que el juez no es lo bastante fiable en fidelidad ni en corrección; se reporta como resultado negativo y no se apoya ninguna conclusión en sus veredictos.

  \item[OE-07 (ejecución de la evaluación y análisis de errores).] \textbf{Conseguido, con la potencia como límite reconocido.} Se ejecutó la evaluación y se analizaron los modos de fallo (el colapso del léxico ante lenguaje ciudadano, la indulgencia del juez, la sobre-abstención del generador). Las conclusiones de recuperación están razonablemente potenciadas; las de generación, apoyadas en un gold revisado —ampliado a las 28 preguntas respondibles de \emph{test}— pero aún de tamaño modesto, se reportan como indicativas.

  \item[OE-08 (reproducibilidad y uso responsable).] \textbf{Conseguido.} El sistema se sostiene íntegramente sobre software de código abierto y modelos de pesos abiertos ejecutados localmente, con artefactos trazables por sus huellas de comprobación, y con un análisis explícito de sus límites técnicos, jurídicos y éticos.
\end{description}

En conjunto, la parte de recuperación queda bien respaldada por los experimentos; la de generación queda establecida y acotada por el tamaño de la anotación y por la insuficiencia del juez. Ninguna conclusión fuerte se apoya en una métrica no validada.

\section{Aplicación de lo aprendido}
\label{sec:aplicacion}

El Trabajo Fin de Grado ha integrado conocimientos de buena parte del Grado en Ciencia e Ingeniería de Datos de la Universidad Rey Juan Carlos. Las asignaturas con una aplicación más directa, agrupadas por área, son las siguientes.

\begin{description}
  \item[Inteligencia artificial, PLN y aprendizaje automático.] \emph{Procesamiento de Lenguaje Natural y Minería de Texto} aportó los fundamentos más directamente aplicados —RAG, tokenización, lematización y \emph{stemming}, \emph{stopwords}, \emph{embeddings} y minería de texto—, presentes en el analizador de español del recuperador léxico, en la fragmentación y en la propia arquitectura RAG. \emph{Aprendizaje Automático I} y \emph{II} proporcionaron la base sobre modelos, representaciones vectoriales, evaluación y la disciplina de separar ajuste (\emph{development}) de prueba (\emph{test}). La asignatura \emph{Generative AI}, cursada en la estancia Erasmus, supuso el primer contacto con los \emph{transformers} y los modelos generativos que sostienen el generador y el juez.

  \item[Ingeniería del software y de datos.] \emph{Fundamentos de Programación} y \emph{Estructuras de Datos y Programación Orientada a Objetos} sustentan la biblioteca modular y el diseño por contratos. \emph{Sistemas Distribuidos de Procesamiento de Datos II} aportó la gestión de proyectos y dependencias (con \texttt{uv}) y las buenas prácticas de ingeniería que se reflejan en el entorno reproducible, la suite de pruebas sin red y el lintado como puerta de calidad. \emph{Diseño y Análisis de Algoritmos} fundamentó la elección de un índice exacto frente a uno aproximado y la fusión de rankings. \emph{Bases de Datos} y \emph{Bases de Datos No Relacionales} inspiran la representación \emph{parent-child} resuelta por \emph{join} y la propiedad única del texto.

  \item[Matemáticas y estadística.] \emph{Álgebra Lineal} está detrás de los \emph{embeddings}, la similitud coseno y el índice denso normalizado. \emph{Inferencia Estadística} y \emph{Probabilidad y Simulación} sostienen el protocolo experimental —remuestreo \emph{bootstrap}, intervalos de confianza y comparación pareada— y los coeficientes de acuerdo ($\kappa$, AC1) con los que se validó el juez.

  \item[Dominio legal, ético y de seguridad.] \emph{Aspectos Legales y Éticos} enmarca el capítulo de impacto y responsabilidad, el tratamiento de los textos del BOE y la protección de datos. \emph{Seguridad y Privacidad de Datos} se refleja en la ejecución local, la no persistencia de las consultas y la gestión de secretos.

  \item[Cómputo y comunicación de resultados.] \emph{Arquitecturas de Altas Prestaciones y Computación en la Nube} enmarca la indexación puntual en GPU, la inferencia en CPU y el análisis de coste computacional. \emph{Visualización de Información} orienta la presentación de tablas y figuras de resultados.
\end{description}

\section{Lecciones aprendidas}
\label{sec:lecciones_aprendidas}

\begin{enumerate}
  \item \textbf{Medir vence a la intuición.} La fusión híbrida «parecía» ayudar en un corpus pequeño y dejó de hacerlo al medirla con rigor sobre 92 normas. Sin el experimento, se habría adoptado una complejidad que no aporta.
  \item \textbf{Un resultado negativo también es un resultado, y a un evaluador automático no hay que creerlo sin validarlo.} La misma cautela que se pide al generador se aplicó al juez: solo tras contrastarlo con anotación humana se supo que no era fiable, y reportarlo con sus datos vale más que exhibir una métrica automática cómoda pero sin respaldo.
  \item \textbf{La distancia entre el lenguaje ciudadano y el normativo es real y medible.} El desplome del léxico en las consultas ciudadanas lo confirma con números, no con impresiones.
  \item \textbf{La dificultad del dominio está en el detalle.} Resolver correctamente la vigencia temporal de cada precepto fue más delicado —y más importante— que cualquier ajuste de modelo.
  \item \textbf{La ingeniería reproducible desde el principio ahorra trabajo.} Los contratos de datos, las configuraciones versionadas, las semillas registradas y una suite de pruebas sin red hicieron posible razonar sobre el sistema y proporcionan los medios necesarios para \emph{regenerar} las tablas del capítulo de experimentos.
\end{enumerate}

\section{Trabajos futuros}
\label{sec:trabajos_futuros}

El prototipo deja abiertas líneas de continuación claras, ordenadas de más inmediatas a más ambiciosas.

\begin{itemize}
  \item \textbf{Reforzar la evaluación de la generación.} Sobre el gold de generación ya ampliado a las 28 preguntas respondibles de \emph{test}, incorporar un segundo anotador que permita medir el acuerdo inter-humano y dar un patrón más sólido de fidelidad.
  \item \textbf{Fidelidad automática sin depender del juez.} Explorar la verificación por implicación textual (\emph{NLI}) como alternativa o complemento al LLM-as-judge, o un juez más potente cuya fiabilidad se valide de nuevo.
  \item \textbf{Reducir la sobre-abstención sin comprometer la seguridad.} El análisis caso a caso de las abstenciones (\ref{subsec:analisis-abstenciones}) descartó el atajo aparente —un umbral sobre la puntuación de recuperación, que separa muy bien lo ajeno al corpus pero no una evidencia suficiente de otra que se le parece \emph{dentro} de él—. Y el experimento con un generador mayor (\ref{subsec:gen-tamano}) mostró que el tamaño, por sí solo, reduce la sobre-abstención pero reintroduce respuestas indebidas fuera del corpus. Las palancas reales son, por tanto, un ensamblado que entregue los párrafos pertinentes con menos ruido y una señal de suficiencia basada en el contenido de la evidencia —no en su puntuación—, que permita responder más sin renunciar a la seguridad.
  \item \textbf{Reordenación y ampliación del corpus.} Evaluar un \emph{reranker} de tipo \emph{cross-encoder} y ampliar el corpus a más áreas normativas, con las técnicas de búsqueda aproximada que la escala requiera.
  \item \textbf{Producto y arquitecturas avanzadas.} Una interfaz web y una API de servicio, y la exploración de arquitecturas RAG más sofisticadas —agénticas o basadas en grafos de conocimiento—, una vez consolidada la base con la que este trabajo se ha comprometido.
\end{itemize}
